%% ==========================================================================
%% SECTION 5: ESCAPING THE IMPOSSIBILITY
%% ==========================================================================
\section{Escaping the Impossibility}\label{sec:escape}

The impossibility results in \autoref{sec:formal} apply to a specific
architectural class: systems operating under text-only observation models with
bounded supervision. They are, in essence, a consequence of unobservability:
the interface discards the internal state that verification requires. They do
not preclude epistemically honest AI systems in general. They characterize the
conditions under which honesty guarantees become achievable, which is to say,
the conditions under which the relevant internal state becomes observable.

We identify three architectural properties that, together, place a system
outside the impossibility class:

\begin{principle}[State Exteriority]
The representation of validity must be separated from the representation of
generation. The policy must explicitly condition on retrieved world-state $w$,
not just query $q$. This breaks the representational impossibility
(\autoref{thm:representational}) by giving the system access to the
information needed to distinguish answerable from unanswerable queries.
\end{principle}

What counts as ``external reality'' matters. A web corpus is external to the
model but not independently verified; it may contain the same fabrications the
model would generate. State Exteriority requires grounding in sources with their
own integrity guarantees: curated databases, sensor data, cryptographically
signed records, or oracles with verifiable provenance.

QuCo-RAG exemplifies partial progress: grounding retrieval decisions in
pre-training corpus statistics rather than model-internal confidence scores. But
corpus co-occurrence is a proxy for truth, not truth itself. And
retrieval-augmented systems can still hallucinate and fabricate
citations because State Exteriority is necessary but not sufficient.

\begin{principle}[Verification Independence]
Verification signals must originate from a channel orthogonal to the generation
signal, capable of overriding the bounded supervisor's preference for
plausibility. This breaks the learnability impossibility
(\autoref{thm:learnability}) by providing gradient signal that distinguishes
honest from fabricated responses.
\end{principle}

Recent work on honesty elicitation demonstrates one implementation: training a
separate verification head whose reward is decoupled from task performance. The
model has no incentive to lie in the verification channel because doing so
doesn't improve its task score. But if the verification head produces text
evaluated by bounded human raters, it remains subject to the same observation
limitations. Full Verification Independence requires the verification channel to
access internal state or external ground truth, not just a second text output.

\begin{principle}[Provenance Binding]
Every output assertion must be structurally bound to a verifiable source. This
addresses the composition problem: as responses grow in complexity, cross-claim
consistency checking becomes superlinear. Provenance binding makes verification
modular because each claim can be checked against its source independently.
\end{principle}

Current retrieval-augmented systems gesture toward this but rarely achieve it.
Citations are generated, not bound. The model produces text that looks like it
cites sources, but the citation is itself a generation subject to fabrication.
True provenance binding requires structural guarantees: signed queries to
retrieval services, content hashes of retrieved passages, or justification
graphs that can be independently audited.

\subsection{Quality of Service Tiers}

These principles suggest a framework for matching system architecture to
epistemic requirements:

\paragraph{Tier 0: Best-Effort (No Guarantees).}
Standard text-only interface, bounded supervision. Suitable for brainstorming,
creative work, low-stakes exploration. The user accepts that any factual claim
may be fabricated. This is not a failure mode, rather it is an appropriate
architecture for tasks where epistemic guarantees are unnecessary.

\paragraph{Tier 1: Bounded Verification (Domain-Constrained).}
Retrieval augmentation with human-in-loop for consequential claims. The claim
space is narrowed to domains where verification is tractable. Explicit
uncertainty surfacing. Suitable for business analysis, civil legal research,
educational applications. The judge is still bounded, but the domain constraint
makes the budget sufficient.

\paragraph{Tier 2: Strong Verification (High-Stakes Domains).}
Narrow claim sets with mandatory provenance chains. External grounding required
for every assertion. Fail-safe to human judgment on any uncertainty. Suitable
for criminal law, medical diagnosis, nuclear safety, aerospace, critical
infrastructure. The cost is high, the domain is severely constrained, but the
guarantees are real.

\medskip

The failure mode is not using AI in high-stakes domains. The failure mode is
claiming Tier~2 guarantees while delivering Tier~0 architecture.

Current frontier systems market themselves as general-purpose assistants
suitable for any task. Our results show this framing is epistemically
incoherent: a system that provides no epistemic guarantees cannot be suitable
for tasks that require epistemic guarantees. Honest marketing would make the
tier explicit.

\subsection{The Path Forward}

The impossibility we prove is not a counsel of despair. It is a map.

Systems engineers have solved analogous problems before. Vector clocks for
causal ordering. Checksums for data integrity. Cryptographic signatures for
authenticity. The tools exist. The architectural patterns exist.

The work ahead is building AI systems that export their epistemic state the way
distributed systems learned to export their causal state. Not as an
afterthought. Not as an optional feature. As a fundamental design requirement.

Concretely: a response under an epistemically-rich interface might include not
just text, but structured metadata, such as answer status (answer, abstain, or
uncertain), provenance (retrieved sources or explicit ``no grounding''),
alternative hypotheses the system considered, and falsifiers (what evidence
would change the conclusion). Such metadata is harder to fabricate than fluent
text, especially when coupled to verifiable external state. The design space is
large; the constraint our theorems impose is that \emph{some} epistemic export
is required because the text-only interface provably cannot suffice.

The interface is the problem. The interface is also where the solution lives.

\subsection{The Tensor Interface: An Existence Proof}

We propose a minimal interface that demonstrates the escape is achievable with
current architectures. The key insight is the distinction between
\emph{testimony} (what the model says) and \emph{telemetry} (measurements of
the computation that produced the output).

\begin{definition}[Tensor Interface]
A generation interface that returns:
\begin{enumerate}
\item \textbf{Text}: The linearized response string (standard interface output).
\item \textbf{Entropy trace}: Per-token entropy of the output distribution during
  generation.
\item \textbf{Attention summary}: Statistics of attention patterns (concentration,
  self-attention ratio) from reasoning layers.
\end{enumerate}
\end{definition}

This interface exports telemetry that the model cannot separately control. The
model can lie in its text, but it cannot present a different forward pass than
the one it actually computed. When the model generates ``Dr. Tanaka's 2023
paper found that...'', its entropy trace reveals whether the generation
proceeded with high confidence (low entropy, suggesting retrieval of stored
pattern) or uncertainty (high entropy, suggesting fabrication).

\paragraph{Implementation.}
The tensor interface requires no architectural changes to current transformers.
Generation APIs already support returning logits; the entropy trace is a simple
transformation. Attention patterns are available with \texttt{output\_attentions=True}
in standard frameworks. The interface is:

\begin{verbatim}
  generate_with_tensor(prompt) ->
    (text, entropy_trace, attention_summary)
\end{verbatim}

We provide a reference implementation supporting OLMo models that demonstrates
the interface is practical. The overhead is minimal: storing per-token entropy
and summarizing attention patterns adds negligible latency compared to generation
itself.

\paragraph{What the tensor provides.}
The entropy trace provides a falsifiability signal. If a model claims certainty
(``Paris is definitely the capital of France'') while its entropy trace shows
high uncertainty, the mismatch is detectable. The attention summary provides
complementary signal: fabrications tend to show higher fragmentation (disconnected
attention patterns) than grounded responses.

\paragraph{What the tensor does not provide.}
The tensor interface is not a lie detector. It does not prove that low-entropy
outputs are true or that high-entropy outputs are false. It provides \emph{one
additional signal} that, combined with other verification methods, can improve
epistemic assessment. The theorems in \autoref{sec:formal} prove that
\emph{some} additional signal is necessary; the tensor interface demonstrates
that such signals are available without architectural change.

\paragraph{Quantitative demonstration.}
To validate the tensor interface under the exact assumptions of our impossibility
theorems, we evaluate bounded verification on 800 query--response pairs (200
queries $\times$ 4 architectures). The full results appear in
\autoref{tab:budget}; the key finding is that tensor-guided selection at 10\%
budget (82.1\% accuracy) outperforms text-guided selection at 30\% budget
(80.4\%). The tensor enables a bounded supervisor to achieve better outcomes
while spending one-third the verification resources, because entropy-based
triage concentrates the budget on the queries most likely to contain
fabrications rather than spreading it uniformly.

\subsection{Tensor-Gated Composition}

The tensor interface generalizes to composed systems. Recent work on Recursive
Language Models (RLMs) demonstrates that LLM outputs can feed into subsequent
LLM calls, enabling unbounded context through recursive
decomposition~\cite{zhang2025recursivelanguagemodels}. However, RLM reports failure modes
characteristic of composition without epistemic state: models abandon correct
intermediate answers, make redundant verification calls, and lose track of what
they have already solved.

The tensor provides a missing primitive: epistemic metadata that can propagate
across composition boundaries. We demonstrate this with tensor-gated
composition, where entropy thresholds control whether outputs propagate to
subsequent calls.

\paragraph{Experimental setup.}
We test three query chains: two fabrication prompts (``What is Glavinsky
syndrome?'' and ``Describe the Brennan-Kowalski theorem'') and one grounded
control (``What is the capital of France?''). For each, we generate an initial
response with tensor extraction, then use an entropy threshold ($\tau = 0.6$)
to gate whether the followup query proceeds.

\paragraph{Results.}
The grounded query produces entropy 0.02, well below threshold; the followup
proceeds and generates accurate content about Parisian landmarks. Both
fabrication queries produce entropy above 0.8; the gate blocks followup queries
that would have requested treatments for the nonexistent syndrome or
implications of the fictional theorem.

This is a minimal demonstration, not a complete solution. But it establishes
that tensor metadata, when used to gate composition, can prevent fabrication
chains from propagating. The supervisor, whether human or automated, gains
information that the text-only interface discards.

\paragraph{Symmetric tensor interface.}
Full compositional integrity requires a symmetric interface: not just
$f(\text{prompt}) \to \text{Tensor}$, but
$f(\text{Tensor}_{\text{in}}) \to \text{Tensor}_{\text{out}}$, where each call
receives the prior call's epistemic context and extends it. For fresh prompts,
the input tensor is empty. For composed calls, the input tensor carries
accumulated provenance and epistemic state. This design enables epistemic
traces to thread through recursive pipelines, providing verification handles at
each composition boundary.

\paragraph{Compositional defense.} The tensor's blind spots characterize where other verification methods are most effective. \autoref{sec:escape}'s three principles — State Exteriority, Verification Independence, Provenance Binding — need not be satisfied by a single mechanism. The tensor provides Verification Independence: an orthogonal signal the model cannot separately control. Bounded judges provide State Exteriority: external lookup against verifiable sources. Input-side defenses (prompt analysis, supply chain detection) provide Provenance Binding at the system boundary. Each layer addresses the failure class where the others are weakest. The tensor catches uncertain fabrication where verification is expensive; the bounded judge catches confident fabrication where verification is cheap; input-side analysis catches adversarial inputs designed to exploit well-learned patterns. Composing these layers allocates a finite verification budget across failure classes rather than spreading it uniformly across the full output space.
